\documentclass[12pt]{article}
\usepackage{amsmath, amssymb}
\usepackage[margin=1in]{geometry}
\setlength{\parskip}{6pt}
\title{QUBO Formulation: HP Lattice Protein Conformation Selection}
\date{}
\begin{document}
\maketitle

\subsection*{HP Lattice Protein Conformation Selection}

\paragraph{Problem Statement.}
Given a hydrophobic-polar (HP) protein sequence and a curated set of $N$ candidate lattice conformations, each conformation $i \in \{0, \dots, N-1\}$ is associated with a contact energy $E_i < 0$ that quantifies the stability contribution of nonconsecutive hydrophobic-hydrophobic contacts. The task is to select exactly one conformation such that the total contact energy is minimized.

\paragraph{Decision Variables.}
For each candidate conformation $i \in \{0, \dots, N-1\}$, introduce a binary decision variable $x_i \in \{0, 1\}$. The variable $x_i = 1$ indicates that conformation $i$ is selected, and $x_i = 0$ otherwise.

\paragraph{QUBO Derivation.}
Step 1 --- Original Objective. We wish to minimize the total contact energy over the binary decision variables:
\[
\min_{x \in \{0,1\}^N} \sum_{i=0}^{N-1} E_i x_i.
\]

Step 2 --- Constraint. Exactly one conformation must be selected, which is expressed as the equality constraint:
\[
\sum_{i=0}^{N-1} x_i = 1.
\]

Step 3 --- Penalty Term Expansion. We enforce the equality constraint by adding a quadratic penalty term to the objective. Let $\lambda > 0$ be the penalty weight. The penalty term is written as:
\[
P(x) = \lambda \left( \sum_{i=0}^{N-1} x_i - 1 \right)^2.
\]
Expanding the square and applying the idempotent property $x_i^2 = x_i$ for binary variables yields:
\begin{align}
P(x) &= \lambda \left( \sum_{i=0}^{N-1} x_i^2 + \sum_{i \neq j} x_i x_j - 2 \sum_{i=0}^{N-1} x_i \right) \\
&= \lambda \left( \sum_{i=0}^{N-1} x_i + 2 \sum_{0 \leq i < j \leq N-1} x_i x_j - 2 \sum_{i=0}^{N-1} x_i \right) \\
&= -\lambda \sum_{i=0}^{N-1} x_i + 2\lambda \sum_{0 \leq i < j \leq N-1} x_i x_j + \lambda.
\end{align}
The resulting general penalty contribution to the Hamiltonian consists of linear terms with coefficient $-\lambda$, quadratic terms with coefficient $2\lambda$, and a constant offset $\lambda$.

Step 4 --- Combined Hamiltonian. The full QUBO Hamiltonian $H(x)$ is the sum of the original objective and the penalty term:
\begin{align}
H(x) &= \sum_{i=0}^{N-1} E_i x_i - \lambda \sum_{i=0}^{N-1} x_i + 2\lambda \sum_{0 \leq i < j \leq N-1} x_i x_j + \lambda \\
&= \sum_{i=0}^{N-1} (E_i - \lambda) x_i + 2\lambda \sum_{0 \leq i < j \leq N-1} x_i x_j + \lambda.
\end{align}

Step 5 --- Q Matrix Entries and Offset. Reading off the coefficients from the general closed form $H(x) = x^\top Q x + c$, we obtain:
\begin{align}
Q_{i,i} &= E_i - \lambda, \quad \text{for all } i \in \{0, \dots, N-1\}, \\
Q_{i,j} &= 2\lambda, \quad \text{for all } 0 \leq i < j \leq N-1, \\
Q_{j,i} &= 2\lambda, \quad \text{for all } 0 \leq i < j \leq N-1, \\
c &= \lambda.
\end{align}

\paragraph{Penalty Weight Justification.}
The penalty weight $\lambda$ must be chosen sufficiently large to ensure that any violation of the equality constraint incurs a higher energy cost than the maximum possible reduction in the objective function. The maximum possible decrease in the objective function by selecting a different conformation is bounded by $\max_{i} |E_i|$. Therefore, setting $\lambda > \max_{i} |E_i|$ guarantees that any infeasible assignment (where $\sum_i x_i \neq 1$) yields a penalty of at least $\lambda$, which strictly dominates any feasible objective improvement. This ensures the global minimum of $H(x)$ over $\{0,1\}^N$ corresponds to a feasible solution.

\paragraph{Correctness Argument.}
Minimizing $H(x)$ over binary vectors recovers the optimal conformation selection because (i) for any feasible solution satisfying $\sum_i x_i = 1$, the penalty term vanishes and $H(x)$ reduces exactly to the original objective $\sum_i E_i x_i + \lambda$; (ii) for any infeasible solution where $\sum_i x_i \neq 1$, the penalty term evaluates to at least $\lambda$. Since $\lambda > \max_i |E_i|$, the energy of any infeasible solution strictly exceeds the energy of the best feasible solution. Consequently, the global minimizer of $H(x)$ must be feasible and, by construction, minimizes the original contact energy.

\paragraph{Illustrative Example.}
Consider a small instance with $N = 3$ candidate conformations and contact energies $E_0 = -5$, $E_1 = -3$, and $E_2 = -4$. We choose $\lambda = 6$, which satisfies $\lambda > \max\{|E_0|, |E_1|, |E_2|\} = 5$.
Substituting these values into the general formulas yields the concrete Hamiltonian:
\begin{align}
H(x) &= (-5 - 6)x_0 + (-3 - 6)x_1 + (-4 - 6)x_2 + 12(x_0 x_1 + x_0 x_2 + x_1 x_2) + 6 \\
&= -11x_0 - 9x_1 - 10x_2 + 12x_0 x_1 + 12x_0 x_2 + 12x_1 x_2 + 6.
\end{align}
The corresponding $Q$ matrix entries and offset are:
\begin{align}
Q_{0,0} &= -11, \quad Q_{1,1} = -9, \quad Q_{2,2} = -10, \\
Q_{0,1} = Q_{1,0} = Q_{0,2} = Q_{2,0} = Q_{1,2} = Q_{2,1} &= 12, \\
c &= 6.
\end{align}
Enumerating all $2^3 = 8$ binary assignments, the feasible solutions are $(1,0,0)$, $(0,1,0)$, and $(0,0,1)$. Their energies are $H(1,0,0) = -11 + 6 = -5$, $H(0,1,0) = -9 + 6 = -3$, and $H(0,0,1) = -10 + 6 = -4$. The infeasible solutions all incur a penalty of at least $6$, yielding energies strictly greater than $-5$. Thus, the global minimum is attained at $x^* = (1,0,0)$ with minimum energy $-5$, correctly selecting conformation $0$.

\end{document}